%! TEX root = pot_temu6.tex
% the main TeX file which is intended to compile, :VimtexReload after adjustment
% :h vimtex-tex-root
\documentclass[../main.tex]{subfiles}
% \includeonlyframes{current}
% \setbeameroption{hide notes} % Only slides
% \setbeameroption{show only notes} % Only notes
% \setbeameroption{show notes on second screen=right} % Both

\begin{document}

\begin{frame}[mytitle=standard,c,mycolor=digiPH_uniblue,light]

	\begin{table}
		\caption{Tugas}\label{tab:}
		\begin{center}
			\begin{tabular}[c]{lp{.4\textwidth}}
				\hline
				Tampak 4 sisi   & : skala 1:200 / 1:250 \\
				Potongan 2 sisi & : skala 1:200 / 1:250 \\
				\hline
			\end{tabular}
		\end{center}
	\end{table}
\end{frame}
\begin{frame}[mytitle=standard,c,mycolor=digiPH_purple,dark] \frametitle{Contoh} \end{frame}
\bgdarkimg{pot1}
\bgdarkimg{pot3}
\bgdarkimg{pot4.png}
\bgdarkimg{pot4.jpg}
\bgdarkimg{tampak1}
\bgdarkimg{tampak2}
\bgdarkimg{tampak3}
\bgdarkimg{tampak4}
\bgdarkimg{tampak5}

\begin{frame}[mytitle=standard,c,mycolor=digiPH_uniblue,light]

	\begin{table}
		\caption{Tugas Lanjutan}\label{tab:}
		\begin{center}
			\begin{tabular}[c]{lp{.4\textwidth}}
				\hline
				Rencana Lanskap & : skala 1:200 / 1:250 \\
				\hline
			\end{tabular}
		\end{center}
	\end{table}
\end{frame}

\begin{frame}[mytitle=standard,c,mycolor=digiPH_purple,dark] \frametitle{Contoh Tugas Lanjutan} \end{frame}

\bgdarkimg{lanskap1}
\bgdarkimg{lanskap2}
\bgdarkimg{lanskap3}
\bgdarkimg{lanskap4}



\end{document}
